\scanpage{196}
\mspage{L 19}
que en l'occurrence) ne s'apparente pas au pur travail du marteau sur l'enclume ou sur le
\unsure{burin} --- le mépris des processus créateurs \inserted{les} plus délicats
(et \inserted{parfois,} de moindre apparence), \deleted{et à la limite,}
de tout ce qui est \deleted{\unsure{idée, vision,}} inspiration, rêve, vision
\deleted{ou \deleted{simplement}\inserted{\deleted{même}} \unsure{idée, même si clairement conçue et formulée soit-elle}}
\inserted{si puissants et \unsure{futiles|fertiles} soient-ils}\inserted{et même}
(à la limite) de toute i\,d\,é\,e, si clairement conçue et formulée soit-elle;
de tout ce qui n'est écrit et p\deleted{\unsure{u}}\,u\,b\,l\,i\,é noir sur blanc,
sous forme d'énoncés purs et durs,\deleted{\unsure{et avec} démontrés}
répertoriables et répertoriés, mûrs pour les ``banques de \uline{données}'' s'engouffrant
dans \deleted{\unsure{d'}}\inserted{les} inépuisables mémoires de nos mégaordinateurs.

\margin{\figure{\begin{tikzpicture}[scale=.3]
\draw[blue] (0,0)--(2,0)--(2,1.5);\draw[blue] (1.8,-.1)--(1.8,1.5);
\end{tikzpicture}}}
Il y a eu \deleted{\unsure{un}} (pour reprendre \deleted{les} une expression de C.L. Siegel
($\unsure{*|x}$)) un extraordinaire ``\uline{aplatissement}'' ou ``\uline{rétrécissement}''
dans la pensée mathématique($\unsure{*|x}$), dépouillée d'une dimension essentielle,
\deleted{d'un versant de son}\inserted{de tout son}
``versant \deleted{flou}\inserted{d'ombre}'' \deleted{ou}\inserted{du versant} ``féminin''.
\deleted{Traditionnellement,}\inserted{Il est vrai que de tradition ancestrale,}
ce versant-là \deleted{restait}\inserted{du travail de\deleted{\unsure{s}} découverte
\deleted{scientifique,} restait} dans une large mesure occulté,
\deleted{\unsure{dans l'inconscient}} personne (autant dire) n'en \uline{parlait} jamais
--- mais le contact avec les sources profondes du rêve, qui alimentent
\deleted{les grands desseins et} les grandes visions et les grands desseins,
\deleted{n'a jamais été}\inserted{n'avait} jamais\inserted{encore} (à ma connaissance)
\unsure{était|été} perdu. Il semblerait \deleted{que}\inserted{qu'à présent}
nous sommes déjà entrés dans une époque
\deleted{où ce contact est sur le point d'être}\inserted{de d\,e\,s\,s\,è\,c\,h\,e\,m\,e\,n\,t, où}
cette source est, non point tarie certes, mais où l'accès à elle est condamné\deleted{\unsure{e}},
par le verdict sans appel \deleted{de la dérision et} du mépris
\deleted{\unsure{général}}\inserted{général} et par les représailles de la dérision.
Nous approchons du moment, semble-t\unsure{è|'}il, où\inserted{sera éradiqué}
non seulement le \uline{souvenir}\inserted{en chacun}
\deleted{\unsure{des stades de la pensée}}
\inserted{de \deleted{ces} tout travail \deleted{\unsure{à notre}} proche de la source,
du travail \inserted{à}}\deleted{\unsure{à la}} dominante ``féminine''
\deleted{dans le travail mathématique (stades}
(ridiculisés par des épithètes comme ``vaseux'', ``mou'', ``inconsistant'' --- ou au bout opposé,
\deleted{``trivial'', ``évident''}\inserted{``trivialités'', ``évidences''}
``bombinage'' \deleted{\unsure{etc}} \ldots), mais
\deleted{des stades}\inserted{\inserted{où sera extirpé,} ce travail}
\deleted{étapes elles-mêmes, celles qui où naissent les}\inserted{de gestation même: celui}
où sont conçues, s'élaborent et naissent les notions et les visions nouvelles.
\inserted{Ce sera}\deleted{le moment}\inserted{\unsure{l'époque}}\inserted{aussi}
où l'exercice de notre art sera réduit à \deleted{\unsure{du|un}}
\inserted{\deleted{pur}\unsure{aride|de}} exercice de ``poids et haltères''\inserted{cérébral}
--- \deleted{celui}\inserted{celle} d'une hypertrophie\inserted{fiévreuse et} stérile,
prenant la suite de plus de trois siècles \deleted{de}\inserted{d'} renouvellement créateur\ldots
\margin{\figure{\begin{tikzpicture}[scale=.3]
\draw (0,3)..controls (-.3,2) and (-.4,1)..(-.4,0);
\end{tikzpicture}}}

Mais à nouveau je\deleted{\unsure{anticipe ou}} digresse, \deleted{et} j'anticipe sur ce que la réflexion m'a enseigné,
\deleted{dans la \unsure{continuation}\deleted{dans la direction}}\inserted{en partant}
d'un double propos clairement présent en moi
\deleted{en commençant mon ``introduction'', en juin 1983 d'abord et en la reprenant en février
l'ann\unsure{ée} dernier (l'année d'après):}
\inserted{dès avant même les débuts de \deleted{\unsure{celle}} celle-ci:}
le propos d'une \inserted{``}déclaration d'intentions\inserted{''}, et (intimement lié à celui-ci)
celui de m'exprimer au sujet de \unsure{la} nature du travail créateur.
Il y avait pourtant un troisième propos encore, moins clairement présent sûrement au niveau conscient,
mais correspondant à un besoin plus profond
\margin{\figure{\begin{tikzpicture}[scale=.3]
\draw (1.8,3.1)--(0,3)--(-.1,2)..controls (-.05,1) and (0,.5)..(.05,0);
\end{tikzpicture}}}
\margin{195}
